\chapter{Conclusion and Perspectives}

\section{Summary of contributions}

This internship delivered a practical adaptive platform for evidence collection on
unstable streaming websites. The work addressed a concrete operational need at Soft Stars:
certain page conditions — deeply nested iframes, heavy anti-bot pressure, click-gated
players, and tokenized stream generation — make fully automated evidence collection
challenging. The agentic platform built during this internship tackles precisely those
cases through reasoning-driven browser automation.

The contributions can be summarized along five dimensions:

\begin{table}[H]
\centering
\begin{tabularx}{\textwidth}{>{\raggedright\arraybackslash\bfseries}p{3.5cm}
                              >{\raggedright\arraybackslash}X}
\toprule
\rowcolor{accent!8}
Contribution & Description \\
\midrule
n8n prototype & Rapid feasibility validation of the extraction pipeline on real targets;
  established that the multi-agent approach was viable before committing to a full coded
  platform. \\
Agent architecture & Four specialist agents (classification, landing, hosting, embedded)
  each with a targeted ReAct-style prompt, dedicated tool set, and explicit stopping
  criteria. \\
Browser tool layer & Puppeteer-based tools (inspect, interact, harvest, screenshot,
  navigate, get\_media\_state) that enable adaptive page interaction beyond static HTML
  retrieval. \\
Operator platform & FastAPI backend with structured observability + Next.js operator
  console providing live run streaming, history, and evaluation surfaces. \\
Prompt engineering & Iterative evolution from intent-based to ReAct-style prompts,
  evaluated across navigation and player activation use cases, with documented improvement
  in recall and recovery rate. \\
\bottomrule
\end{tabularx}
\caption{Summary of main contributions by dimension.}
\label{tab:contrib-summary}
\end{table}

\section{Key technical lessons}

\subsection{Prompt design matters more than agent count}

The biggest improvement across all evaluation runs came from switching to ReAct-style prompts, not from adding more agents or tools. The intent-based prompt left agents free to skip navigational controls entirely — they'd look at the first visible state and decide they were done. Adding an explicit THINK step that checks for tabs or checks player state before harvesting made the agents dramatically more reliable. More tools without better prompts just gave the agent more ways to make the same mistakes.

\subsection{Fallback paths are not optional}

No single extraction strategy works on all pages. The hosting-to-embedded fallback recovered cases that would have been hard failures — the 3-level iframe chain in Case A being the clearest example. Without that routing rule, the hosting agent would have reported no stream found and the run would have ended there. In this domain, you need layers.

\subsection{Observability drives improvement}

Without the per-step tool trace and screenshot artifacts, most of the debugging during evaluation would have been guesswork. The improvements to tool descriptions — adding the "use only after activation confirmed" guard to harvest, clarifying when to use inspect\_hosting versus inspect\_page — were only discovered by looking at the actual tool call logs, not by reading the prompt. Structured observability is what made it possible to turn failures into targeted fixes.

\subsection{Anti-bot adaptation is continuous}

Nothing static stayed reliable across the evaluation window. Adblock filter lists needed updating. Fingerprint settings that worked in week one triggered challenges in week three on the same sites. The value of the platform is not that it "solves" anti-bot — it doesn't — but that it gives you a framework for continuously re-solving it as things change.

\section{Conclusion}

The platform built during this internship works, but more importantly it's maintainable. That distinction matters in a domain where target behavior changes continuously. A system that works once and then breaks the next time a site updates its anti-bot configuration is not useful for legal enforcement. What Soft Stars needed was a framework they could keep adapting — and that's what the evaluation campaign actually validated.

The improvements didn't come from clever algorithms. They came from systematically working through failure modes: navigation-heavy sites needed state-loop exploration rather than single-pass scanning; popup-blocked pages needed fallback click strategies; Cloudflare-gated sessions needed lower automation signals; tokenized streams needed deferred harvest after play activation. Each fix was small on its own but the combination produced a system that handles the majority of real-world streaming infrastructure effectively.

The n8n phase was genuinely useful as a rapid validation tool. The coded platform replaced it not because n8n was wrong but because the engineering requirements outgrew what a visual workflow tool can sensibly maintain. Versioned prompts, structured traces, per-run cost tracking, and a proper evaluation harness — those are what make it possible to keep improving the system after the internship is done.

\section{Future work}

\subsection{Short-term improvements}

\begin{itemize}
  \item formalize periodic fingerprint and adblock policy rotation as a scheduled
    maintenance process;
  \item expand case-study benchmarks with standardized before/after mitigation evidence;
  \item integrate stricter automatic fallback triggers between deterministic and agentic
    paths;
  \item add a \texttt{scroll-and-reinspect} loop to the landing agent for sites that
    require scrolling to reveal all match cards.
\end{itemize}

\subsection{Medium-term directions}

\begin{itemize}
  \item validate cloud execution under production load profiles (Azure Service Bus +
    Container Apps Jobs);
  \item implement long-term vector memory for site layout pattern reuse across runs;
  \item add YOLO or OCR-based play button detection as a fallback for non-standard player
    UIs;
  \item implement multi-server parallel harvesting to reduce time-per-host-page;
  \item build a prompt regression test suite that runs against a fixed snapshot dataset
    whenever prompts are modified.
\end{itemize}

\subsection{Research directions}

\begin{itemize}
  \item investigate self-improving selector generation: the agent proposes a new selector,
    the system tests it, and stores the result in long-term memory;
  \item explore fine-tuned classification models as replacements for the LLM
    classification call to reduce cost at scale;
  \item study the stability of tokenized stream URL patterns to determine if
    pre-activation network monitoring can predict manifest timing.
\end{itemize}

\section{Deployment roadmap}

\begin{table}[H]
\centering
\begin{tabularx}{\textwidth}{>{\raggedright\arraybackslash\bfseries}p{2.5cm}
                              >{\raggedright\arraybackslash}X
                              >{\raggedright\arraybackslash}X}
\toprule
\rowcolor{accent!8}
Phase & Objective & Azure components \\
\midrule
Phase 1 (done) & Local and containerized execution; cloud-compatible design &
  Docker containers; PostgreSQL \\
Phase 2 (next) & Queue-driven execution on cloud & Azure Service Bus; Container Apps
  Jobs; Blob storage \\
Phase 3 (future) & Auto-scaling and parallel job execution & Container Apps with
  KEDA scaling; multi-replica job workers \\
\bottomrule
\end{tabularx}
\caption{Deployment roadmap toward fully cloud-native execution.}
\label{tab:deploy-roadmap}
\end{table}

\section{Reflection on the methodology}

Using CRISP-DM as an organising framework for an agentic engineering project was a deliberate
choice, but it required adaptation. The standard CRISP-DM phases presuppose a dataset that
can be explored statistically before modelling begins. In this project, the ``data'' is not
a file: it is a live web page whose structure varies at runtime. Data understanding was
therefore conducted through manual observation campaigns against real targets, and the
modelling phase consisted of designing agent prompts and tool schemas rather than fitting
predictive models.

The iterative character of CRISP-DM was, however, precisely the right framework for
prompt engineering. The inner cycle of \textit{observe failure $\to$ hypothesize cause $\to$
revise prompt $\to$ re-test} mirrors the CRISP-DM loop between modeling and evaluation.
Each prompt version recorded in the evaluation chapter corresponds to one CRISP-DM cycle.
The framework gave the iterative process structure and made it easier to document decision
points that would otherwise have been invisible.

A key methodological insight from the evaluation is that \textbf{tool descriptions are
a first-class engineering artifact}. They are not documentation for human readers; they
are the primary signal the LLM uses to decide which tool to invoke on each reasoning
step. Treating tool descriptions with the same care as prompt system messages, and
evaluating them empirically through misuse rate tables, was one of the most effective
levers for improving system reliability without changing the model or the tool
implementation.

\section{Lessons for future agentic projects}

A few things from this project feel worth recording for anyone building something similar.

Specialization helps more than it might seem. The four-agent design costs more to set up than a single general agent, but it pays back quickly: smaller tool menus mean fewer wrong-tool selection errors, prompts are easier to reason about because each one only handles one page type, and failures are easier to trace because you know which agent was responsible. A general-purpose agent that handles landing pages, host pages, and embedded players all at once accumulates context fast and tends to confuse itself.

ReAct-style prompts are not just a style choice — they produce measurably different behavior. The difference is not subtle. An intent-based prompt ("find the stream") leaves the agent free to declare success after looking at the initial HTML. An explicit instruction ("call get\_media\_state before harvest; if idle, click play; if click opens a new tab, close it and retry with js\_click") eliminates entire failure categories. The specificity feels excessive until you watch an intent-based agent harvest an empty page three times in a row.

Observability designed upfront is worth more than observability added later. The decision to log every tool call, screenshot, and LLM response from the start — rather than adding logging when something went wrong — is why the evaluation was productive. Most of the improvements to tool descriptions came from looking at the trace for a failed run and seeing exactly which tool was called in the wrong order or with the wrong timing.

Memory is also what separates a demo from a production tool. The n8n prototype started every run from scratch. The coded platform stores what worked on each domain and reuses it. That cross-run knowledge accumulation is how the system gets faster and more reliable over time without manual intervention.

\section{Final remark}

The final value of the project is not that it ``solves'' anti-bot constraints once.
Its value is that it provides a disciplined framework to keep solving them as target
behavior changes, while preserving technical evidence that investigators can trust.

The clearest sign of success is not the extraction rate alone but the fact that the
platform produces reviewable, auditable evidence with traceable reasoning at every step.
That traceability is what makes the agentic workflow operationally useful for a legal
and investigative context like beIN SPORTS evidence collection.

Beyond the specific use case, this project illustrates a general pattern for building
adaptive automation in adversarial environments: combine a reasoning layer that decides
what to do next with a browser tool layer that executes those decisions, structure the
evidence trail so that every failure is diagnosable, and iterate on prompts and tool
descriptions as first-class engineering artifacts. That pattern is not unique to sports
rights enforcement — it applies to any domain where the target environment is dynamic,
heterogeneous, and deliberately resistant to automation.
